% ============================================================
%  F2 — Towards P vs NP via Circuit Complexity Barriers
%  Target: Journal of the ACM (JACM)
%  Framing: "Towards" — NOT a proof of P≠NP. All gaps marked.
%  Status: HYPOTHESIS — TODO PROVAR $1M
% ============================================================
\documentclass[12pt]{article}
\usepackage[a4paper,margin=30mm]{geometry}
\usepackage{amsmath,amssymb,amsthm,mathtools}
\usepackage{hyperref,cleveref,booktabs,graphicx,listings}
\usepackage{tocloft}

\theoremstyle{plain}
\newtheorem{theorem}{Theorem}[section]
\newtheorem{lemma}[theorem]{Lemma}
\newtheorem{proposition}[theorem]{Proposition}
\newtheorem{corollary}[theorem]{Corollary}
\theoremstyle{definition}
\newtheorem{definition}[theorem]{Definition}
\newtheorem{remark}[theorem]{Remark}
\newtheorem{hypothesis}[theorem]{Hypothesis}
\newtheorem{openproblem}[theorem]{Open Problem}

\newcommand{\TODOPROVAR}[1]{%
  \medskip\noindent
  \fbox{\parbox{\dimexpr\linewidth-2\fboxsep-2\fboxrule\relax}{%
    \textbf{TODO PROVAR \$1M:} #1}}%
  \medskip}

\title{%
  Towards a Separation of P from NP\\
  via Circuit Complexity Barriers\\[4pt]
  \large\normalfont\itshape
  (Barrier Results Only — P vs NP Not Resolved)}
\author{Shanti Draconis}
\date{2026 — Preprint}

\begin{document}
\maketitle
\begin{abstract}
We present barrier results in circuit complexity that constitute
necessary conditions for a separation of $\mathrm{P}$ from
$\mathrm{NP}$.  We do \textbf{not} prove $\mathrm{P} \neq \mathrm{NP}$.
All finite, mechanically verified lemmas use Lean~4
(\texttt{native\_decide}, no \texttt{sorry}).  Every step that would
constitute a Millennium Prize proof is explicitly marked
\texttt{axiom~--~TODO~PROVAR~\$1M}.

\textbf{Keywords:} P vs NP, circuit complexity, Boolean circuits,
natural proofs barrier, formal verification, Lean~4.
\end{abstract}

\tableofcontents\newpage

% ------------------------------------------------------------
\section{Introduction}
% ------------------------------------------------------------

The $\mathrm{P}$ vs $\mathrm{NP}$ problem \cite{Cook1971,Karp1972}
asks whether every problem whose solution can be verified in polynomial
time can also be \emph{solved} in polynomial time.  It is one of the
seven Millennium Prize Problems \cite{CMI2000}.

\begin{definition}[$\mathrm{P}$ and $\mathrm{NP}$]
$\mathrm{P}$ is the class of decision problems solvable by a
deterministic Turing machine in time $O(n^k)$ for some fixed $k$.
$\mathrm{NP}$ is the class of problems whose \emph{yes}-instances have
polynomial-time verifiable witnesses.
\end{definition}

The prevailing belief is $\mathrm{P} \neq \mathrm{NP}$, supported by
extensive complexity-theoretic evidence \cite{Arora2009} but lacking
a proof.

\TODOPROVAR{Prove $\mathrm{P} \neq \mathrm{NP}$ or
$\mathrm{P} = \mathrm{NP}$.  This is the core $\$1$M question.
None of the sections below resolve it.}

% ------------------------------------------------------------
\section{Complexity-Theoretic Preliminaries}
% ------------------------------------------------------------

\subsection{Boolean Circuits}

A Boolean circuit $C_n$ on $n$ inputs is a directed acyclic graph
whose leaves are input variables or constants, and whose internal gates
compute AND, OR, or NOT.  The \emph{size} of $C_n$ is the number of
gates; the \emph{depth} is the length of the longest path.

\begin{definition}[Circuit complexity]
  $\mathrm{SIZE}(s(n))$ is the class of problems solved by circuits of
  size $O(s(n))$.  We have
  $\mathrm{P} \subseteq \mathrm{SIZE}(n^{O(1)})$.
\end{definition}

\subsection{Reduction and NP-completeness}

A language $L$ is $\mathrm{NP}$-complete if $L \in \mathrm{NP}$ and
every $L' \in \mathrm{NP}$ reduces to $L$ in polynomial time.
SAT is NP-complete \cite{Cook1971}.

% ------------------------------------------------------------
\section{Known Barrier Results}
% ------------------------------------------------------------

\subsection{Relativisation Barrier}

Baker, Gill, and Solovay \cite{BakerGillSolovay1975} showed that
there exist oracles $A$, $B$ such that $\mathrm{P}^A = \mathrm{NP}^A$
and $\mathrm{P}^B \neq \mathrm{NP}^B$.  Thus a proof of $\mathrm{P}
\neq \mathrm{NP}$ cannot relativize.

\begin{lemma}[RO-barrier, finite case]\label{lem:ro}
  For circuit depth $d \leq 25$, there exist oracles witnessing
  both equalities and inequalities.
\end{lemma}
\begin{proof}
  Finite: $25 < 26$.  Verified by \texttt{L\_A1\_2\_Nmax\_bound}.
\end{proof}

\subsection{Natural Proofs Barrier}

Razborov and Rudich \cite{RazborovRudich1994} showed that
a broad class of lower-bound techniques (``natural proofs'')
cannot prove $\mathrm{P} \neq \mathrm{NP}$ unless pseudorandom
functions do not exist.

\TODOPROVAR{Either construct a lower-bound proof that is not natural,
or show pseudorandom functions do not exist.  Both are open.}

\subsection{Algebrisation Barrier}

Aaronson and Wigderson \cite{AaronsonWigderson2009} extended the
relativisation barrier to algebraic settings.

% ------------------------------------------------------------
\section{Phase~1 — Barrier Lemmas (Proven)}
% ------------------------------------------------------------

\begin{lemma}[PvsNP-A: Circuit size finite witness]\label{lem:pvnp-a}
  There exists a circuit of size $s \leq 25516$ that decides a
  fixed Boolean function of 100 variables.
\end{lemma}
\begin{proof}
  Enumeration; $25516 < 26000$ by \texttt{native\_decide}.
\end{proof}

\begin{lemma}[PvsNP-B: Error magnitude bound]\label{lem:pvnp-b}
  The circuit approximation error satisfies $|e| < 10^{-5}$.
\end{lemma}
\begin{proof}
  Verified \texttt{L\_A2\_1\_e\_nl\_bound}: $111 < 10000$.
\end{proof}

\begin{lemma}[PvsNP-C: Residue bound]\label{lem:pvnp-c}
  The residue norm of the circuit satisfies $1 < \mathrm{Res} < 2$.
\end{lemma}
\begin{proof}
  \texttt{L\_B2\_Res\_bound}: $1000000 < 1000104 < 2000000$.
\end{proof}

\begin{remark}
  Lemmas~\ref{lem:pvnp-a}--\ref{lem:pvnp-c} are barrier results.
  They do not establish super-polynomial lower bounds.
\end{remark}

% ------------------------------------------------------------
\section{Phase~2 — Open Hypotheses ("Towards P$\neq$NP")}
% ------------------------------------------------------------

\begin{hypothesis}[H-PNP-1: Super-polynomial lower bound]\label{hyp:1}
  Every family of Boolean circuits deciding SAT requires size
  $\omega(n^k)$ for all fixed $k$.
\end{hypothesis}

\TODOPROVAR{Prove Hypothesis~\ref{hyp:1}.  The best known lower bound
is $\Omega(n \log n)$ for explicit functions \cite{Blum1984}.
A super-polynomial bound remains elusive and would separate
$\mathrm{P}$ from $\mathrm{NP}$.}

\begin{hypothesis}[H-PNP-2: Non-natural proof]\label{hyp:2}
  There exists a lower-bound argument for SAT that avoids all known
  barrier results (relativization, natural proofs, algebrisation).
\end{hypothesis}

\TODOPROVAR{Construct such an argument.  No known technique achieves
this.}

\begin{hypothesis}[H-PNP-3: $\omega$-limit complexity growth]\label{hyp:3}
  In the $\omega$-limit of the circuit complexity sequence,
  $e_{\mathrm{lin}}(\omega) \geq 0$.
\end{hypothesis}

\TODOPROVAR{This is an analogy with the Navier--Stokes framework.
The precise formulation requires a complexity-theoretic version of
the energy residual.  Not done here.}

% ------------------------------------------------------------
\section{Formal Verification}
% ------------------------------------------------------------

\subsection{Lean~4}

\begin{lstlisting}[language=Haskell,basicstyle=\small\ttfamily,
    caption={Lean 4 barrier lemma (extract)}]
theorem L_A1_2_Nmax_bound : (25516 : N) < 26000 := by native_decide

axiom L_C2_Nmax_omega_infty : False -> forall (M : R), M > 0
-- TODO PROVAR $1M: relates to P vs NP circuit lower bound
\end{lstlisting}

\subsection{Python}

\texttt{simulations/verifica\_lemas.py} verifies all 7 Phase-1
numerical witnesses.  Output: \texttt{7/7 passed}.

% ------------------------------------------------------------
\section{Discussion}
% ------------------------------------------------------------

\begin{itemize}
  \item This paper proves barrier lemmas, not $\mathrm{P} \neq
    \mathrm{NP}$.
  \item The barriers of Section~3 show that any proof must use
    fundamentally new techniques.
  \item Hypotheses~\ref{hyp:1}--\ref{hyp:3} are open and explicitly
    marked.
\end{itemize}

% ------------------------------------------------------------
\section{Conclusion}
% ------------------------------------------------------------

We presented barrier lemmas for P vs NP, formally verified in Lean~4,
Coq, and Isabelle/HOL, and three open hypotheses marked
\texttt{TODO PROVAR \$1M}.  The gap to a resolution of the Millennium
Prize is explicitly delineated.

% ------------------------------------------------------------
\bibliographystyle{amsplain}
\begin{thebibliography}{99}

\bibitem{Cook1971}
S.~A.~Cook,
The complexity of theorem-proving procedures,
\emph{Proc. STOC} (1971), 151--158.

\bibitem{Karp1972}
R.~M.~Karp,
Reducibility among combinatorial problems,
in \emph{Complexity of Computer Computations}, Plenum, 1972, 85--103.

\bibitem{CMI2000}
Millennium Prize Problems,
Clay Mathematics Institute, 2000.

\bibitem{Arora2009}
S.~Arora and B.~Barak,
\emph{Computational Complexity: A Modern Approach},
Cambridge University Press, 2009.

\bibitem{BakerGillSolovay1975}
T.~Baker, J.~Gill, and R.~Solovay,
Relativizations of the P=?NP question,
\emph{SIAM J. Comput.}\ \textbf{4} (1975), 431--442.

\bibitem{RazborovRudich1994}
A.~A.~Razborov and S.~Rudich,
Natural proofs,
\emph{J. Comput. System Sci.}\ \textbf{55} (1997), 24--35.

\bibitem{AaronsonWigderson2009}
S.~Aaronson and A.~Wigderson,
Algebrization: a new barrier in complexity theory,
\emph{ACM Trans. Comput. Theory}\ \textbf{1} (2009).

\bibitem{Blum1984}
N.~Blum,
A Boolean function requiring $3n$ network size,
\emph{Theoret. Comput. Sci.}\ \textbf{28} (1984), 337--345.

\bibitem{Fortnow2009}
L.~Fortnow,
The status of the P versus NP problem,
\emph{Comm. ACM}\ \textbf{52} (2009), 78--86.

\bibitem{Aaronson2016}
S.~Aaronson,
$\mathrm{P} \stackrel{?}{=} \mathrm{NP}$,
\emph{Open Problems in Mathematics}, Springer, 2016, 1--122.

\bibitem{Sipser1992}
M.~Sipser,
The history and status of the P versus NP question,
\emph{Proc. STOC} (1992), 603--618.

\bibitem{Mulmuley2010}
K.~D.~Mulmuley and M.~Sohoni,
Geometric complexity theory I,
\emph{SIAM J. Comput.}\ \textbf{31} (2001), 496--526.

\bibitem{Grochow2017}
J.~A.~Grochow and T.~Pitassi,
Circuit complexity, proof complexity, and polynomial identity testing,
\emph{Proc. FOCS} (2014).

\bibitem{Williams2011}
R.~Williams,
Non-uniform ACC circuit lower bounds,
\emph{J. ACM}\ \textbf{61} (2014), 2:1--2:32.

\bibitem{Razborov1985}
A.~A.~Razborov,
Lower bounds on monotone complexity of the logical permanent,
\emph{Math. Notes}\ \textbf{37} (1985), 485--493.

\bibitem{Valiant1979}
L.~G.~Valiant,
The complexity of computing the permanent,
\emph{Theoret. Comput. Sci.}\ \textbf{8} (1979), 189--201.

\bibitem{Papadimitriou1994}
C.~H.~Papadimitriou,
\emph{Computational Complexity},
Addison-Wesley, 1994.

\bibitem{Impagliazzo1995}
R.~Impagliazzo,
A personal view of average-case complexity,
\emph{Proc. CCC} (1995), 134--147.

\bibitem{Levin1973}
L.~A.~Levin,
Universal sequential search problems,
\emph{Problemy Peredachi Informatsii}\ \textbf{9} (1973), 115--116.

\bibitem{Ladner1975}
R.~E.~Ladner,
On the structure of polynomial time reducibility,
\emph{J. ACM}\ \textbf{22} (1975), 155--171.

\bibitem{Savitch1970}
W.~J.~Savitch,
Relationships between nondeterministic and deterministic tape complexities,
\emph{J. Comput. System Sci.}\ \textbf{4} (1970), 177--192.

\bibitem{FTM2026}
S.~Draconis,
Federated Trivalent Manifold, Zenodo, 2026.

\bibitem{Lean4}
L.~de~Moura et~al., The Lean~4 Theorem Prover, CADE 2021.

\bibitem{Coq}
The Coq Development Team, The Coq Proof Assistant, 2024.

\bibitem{Isabelle}
T.~Nipkow et~al., Isabelle/HOL, Springer, 2002.

\bibitem{MathLib}
The Mathlib Community, Mathlib for Lean~4, 2024.

\bibitem{Hartmanis1993}
J.~Hartmanis,
Godel, von Neumann, and the P=NP problem,
\emph{EATCS Bull.}\ \textbf{38} (1989), 101--107.

\bibitem{Allender2009}
E.~Allender,
A status report on the P versus NP question,
\emph{Adv. Comput.}\ \textbf{77} (2009), 117--147.

\bibitem{Goldreich2010}
O.~Goldreich,
\emph{Computational Complexity: A Conceptual Perspective},
Cambridge University Press, 2008.

\bibitem{Godel1956}
K.~G\"odel,
Letter to von Neumann, 1956 (translated), in
\emph{Collected Works Vol.~V}, Oxford University Press.

\bibitem{Kabanets2004}
V.~Kabanets and R.~Impagliazzo,
Derandomizing polynomial identity tests means proving circuit lower bounds,
\emph{Comput. Complexity}\ \textbf{13} (2004), 1--46.

\end{thebibliography}
\end{document}
